% ---
% title: MiniMax Loss
% date: 2025/11/3
% tag: ML
% ---
“\textbf{MiniMax 损失函数}”通常指的是 \textbf{极小极大优化目标}，最常见出现在 \textbf{GAN（生成对抗网络）} 里。

\hrule

\section{基本含义}

MiniMax 就是：

\begin{itemize}
  \item 一个模型想让目标函数\textbf{尽可能大}
  \item 另一个模型想让目标函数\textbf{尽可能小}
\end{itemize}

写成数学形式就是：\(\min_G \max_D V(D,G)\)

其中：

\begin{itemize}
  \item \(G\) 是生成器（Generator）
  \item \(D\) 是判别器（Discriminator）
  \item \(V(D,G)\) 是两者博弈的目标函数
\end{itemize}

这就是“极小极大”损失。

\hrule

\section{GAN 中最经典的 MiniMax 损失}

原始 GAN 的目标函数是：

\[
\min_G \max_D V(D,G) = \mathbb E_{x\sim p_{\text{data}}(x)}[\log D(x)] + \mathbb E_{z\sim p_z(z)}[\log(1-D(G(z)))]
\]

含义如下：

\subsection{判别器\textbf{\(D\)}的目标}

希望：

\begin{itemize}
  \item 对真实样本 \(x，D(x)\to 1\)
  \item 对生成样本 \(G(z)，D(G(z))\to 0\)
\end{itemize}

所以它要\textbf{最大化}

\[
\mathbb E_{x\sim p_{\text{data}}}[\log D(x)] + \mathbb E_{z\sim p_z}[\log(1-D(G(z)))]
\]

\subsection{生成器\textbf{\(G\)}的目标}

希望生成的样本骗过判别器，让 \(D(G(z))\to 1\)

但在原始 minimax 写法里，生成器是去\textbf{最小化}

\[
\mathbb E_{z\sim p_z}[\log(1-D(G(z)))]
\]

\hrule

\section{为什么叫损失函数}

因为训练时通常会把上面的目标改写成损失形式。

\subsection{判别器损失}

常写为：

\[
L_D = -\mathbb E_{x\sim p_{\text{data}}}[\log D(x)] -\mathbb E_{z\sim p_z}[\log(1-D(G(z)))]
\]

因为实际训练里一般是“最小化损失”，所以把“最大化目标”前面加个负号。

\subsection{生成器的 minimax 损失}

\[
L_G^{\text{minimax}} = \mathbb E_{z\sim p_z}[\log(1-D(G(z)))]
\]

\hrule

\section{原始 MiniMax 损失的问题}

这个生成器损失虽然理论上对，但训练中常有一个问题：

当判别器很强时，\(D(G(z))\approx 0\)，这时\(\log(1-D(G(z)))\)

对生成器提供的梯度会很弱，容易出现\textbf{梯度消失}。

所以实际中常不用这个原始 minimax 形式训练生成器。

\hrule

\section{更常用的 Non-saturating loss}

实际训练 GAN 时，生成器通常改成最大化 \(\log D(G(z))\)，等价于最小化：

\[
L_G^{\text{NS}} = -\mathbb E_{z\sim p_z}[\log D(G(z))]
\]

这个版本梯度更强，更容易训练。

所以你如果在论文里看到：

\begin{itemize}
  \item \textbf{minimax loss}：多半指原始 GAN 论文中的博弈目标
  \item \textbf{generator loss}：实际实现里往往改成 \(\textbf{non-saturating loss}\)
\end{itemize}

\hrule

\section{从博弈角度理解}

MiniMax 的本质是一个二人零和博弈：

\begin{itemize}
  \item 判别器努力区分真假
  \item 生成器努力伪造得更真实
\end{itemize}

两者不断对抗，最终理想状态下：

\[
p_g = p_{\text{data}}
\]

也就是生成分布接近真实分布。

此时判别器无法区分真假，有：

\[
D(x)=\frac12
\]

\hrule

\section{如果不是 GAN，MiniMax 也可以泛指什么}

更一般地，MiniMax 损失函数就是这种形式：\(\min_\theta \max_\phi f(\theta,\phi)\)

常见于：

\begin{itemize}
  \item 生成对抗网络 GAN
  \item 对抗训练
  \item 域适应
  \item 鲁棒优化
  \item 博弈论建模
\end{itemize}

也就是说，“MiniMax”不是某一个固定公式，而是一类\textbf{优化结构}。
